% Page budget: 1.55 pages | Owns: augmentation topology, additional states, encoder, and closed-loop objective.
% WRITING GUIDE
% Purpose: Define the final augmentation architecture and explain how it is trained inside the known feedback loop.
% Start here: Current implementation decisions D-129, D-140, D-141, D-180, D-220, D-233, and D-237 in docs/decisions.md.
% Supporting material: Quinten's 18 September outline feedback, the Meeting-audit synthesis, Thesis-writeup/Documentation/TRAINING-DESIGN.md, and docs/writeup figures.
% Mandatory papers: Read Hoekstra's 2025 LFR augmentation paper, the 2026 first-principles augmentation paper, the encoder-initialization paper, and Kessels' Chapter 5 before writing this section.
% Research-plan anchor: Preserve Aspect 2's dynamic parallel augmentation objective; document the departure from open-loop training and earlier state routing.
% Current decisions: Separate physical and additional states, keep the controller outside the model, score the complete training window without burn-in, and use one network for the physical correction and added-state update.
% Choices to justify: Final reported routing, added-state dimension, ANN width and depth, zero initialization, encoder history, closed-loop objective, rollout length, full-window scoring, and validation horizon. Use the configuration of the reported thesis run, not a superseded routing experiment.
% Horizon check: Treat encoder history, scored training window, joint-estimation horizon, and free-run validation horizon separately; relate candidates to relevant stable modes and test sensitivity to slower dynamics and free-integrator accumulation.
% Implementation audit: Read scripts/gantry/gantry_interconnect_dynamic.py, gantry_dynamic/{model,config,controller,data,training}.py, and model_augmentation/fit_systems/{interconnect,closed_loop,pre_encoder}.py.
% Reference check: Verify sources for parallel LFR augmentation, encoder initialization, closed-loop state extension, and any claimed architectural precedent against the original paper or thesis.
% Cautions: Earlier routing plans are superseded; distinguish the truth-only hidden absorber from learned states; closed-loop fit alone does not prove autonomous plant recovery.
% Figures: Absorber schematic, author's updated architecture, and thesis-owned common-reference loops. No residual panel or horizon figure.
% Section is finished when: Every signal has a defined frame and timing, the RK4 boundary is explicit, and the described architecture matches the final code.
%
% SOURCE MAP
% Hoekstra et al. 2025 EJC: dynamic-parallel topology, additional states, encoder-based truncated simulation.
% Hoekstra et al. 2026 Automatica submission: extended LFR formulation, baseline-preserving model/encoder initialization, joint identification context.
% Hoekstra et al. 2026 encoder paper: reconstructability-based W^b initialization; it does not identify W^a or test dynamic augmentation.
% Kessels 2025 thesis: direct closed-loop rollout, controller equations, and reconstruction of the machine controller state for each window.
% Drenth 2025 thesis: LPV-specific augmentation precedent; not the unrestricted MLP structure used here.
% This project: gantry specialization, residual-loop implementation, exact routing, RK4 boundary, and verification.
\section{Dynamic LPV-LFR Augmentation}\label{sec:augmentation}

The baseline \eqref{eq:eom} omits dynamics of the machine.
Garc\'ia-Herreros et al.\ neglect the resonance of the supporting base, the
vibration of the cross-arm on its flexible plates, the detent forces of the
linear motors, and the variation of friction along the
stroke~\cite[Sec.~2.4]{garcia2013model}.
Our baseline also leaves out the Coulomb friction that their model contains.
We augment the baseline with a learned dynamic component and estimate its
parameters jointly with $\theta_{\mathrm{base}}$ from records measured under
feedback.
Section~\ref{sec:aug_structure} defines the augmented model,
Section~\ref{sec:aug_encoder} the initial state of each training window, and
Section~\ref{sec:aug_closed_loop} the closed-loop rollout and the objective.

\subsection{Augmented model}\label{sec:aug_structure}
An omitted vibration mode requires additional model order.
Neither a refit of $\theta_{\mathrm{base}}$ nor a static correction of the
baseline state update provides it, because both act only on the existing
states~\cite[Sec.~2]{hoekstra2025lfr}.
We therefore add states and identify the dynamic-parallel
model~\cite[Table~1]{hoekstra2026lfr}
\begin{equation}
\begin{aligned}
 \tilde x_{k+1}
 &=f_{\mathrm{base}}(\theta_{\mathrm{base}},\tilde x_k,u_k)
   +f_{\mathrm{aug}}(\theta_{\mathrm{aug}},\hat x_k,u_k),\\
 \bar x_{k+1}
 &=g_{\mathrm{aug}}(\theta_{\mathrm{aug}},\hat x_k,u_k),\\
 \hat y_k&=h_{\mathrm{base}}(\tilde x_k),
\end{aligned}
\label{eq:aug_transition}
\end{equation}
where $\tilde x=\col(q,\dot q)\in\R^6$ is the physical state of
Section~\ref{sec:lfr}, $\bar x\in\R^{n_{x_a}}$ collects the additional
states, $\hat x=\col(\tilde x,\bar x)$, $u=P\uact$ is the generalised force
of \eqref{eq:Ptransform}, and $\hat y$ is the model estimate of the measured
positions $y=\qact$.

Under joint estimation, the split between $f_{\mathrm{base}}$ and
$f_{\mathrm{aug}}$ is not unique: $f_{\mathrm{aug}}$ can learn changes that a
different $\theta_{\mathrm{base}}$ would
produce~\cite[Sec.~5.2]{hoekstra2026lfr}.
The augmentation is parallel rather than in series because
Section~\ref{sec:obc} constrains this split with an orthogonal construction
that Gy\"or\"ok et al.\ formulate for an additive learning
component~\cite[Sec.~2]{gyorok2025l4dc}.

Drenth augments a self-scheduled discrete-time LPV-LFR baseline with
additional states through extra LFR blocks and scheduling
maps~\cite[Sec.~5.2]{drenth2025thesis}.
Here, the learned part is an unstructured network acting on a
continuous-time physical baseline, so the dynamic-parallel form of
\eqref{eq:aug_transition} is used instead.

$f_{\mathrm{base}}$ is one RK4 step of \eqref{eq:eom} over $T_s$, with the
input held over the step.
Each stage evaluates $M(Y)$ at its own state, so $f_{\mathrm{base}}$ remains
self-scheduled.
All signals are normalised with the statistics of the training
records~\cite[Sec.~5.3]{hoekstra2026lfr}.

One multilayer perceptron (MLP) $\phi_{\mathrm{aug}}$ with tanh hidden layers
supplies both learned functions,
\begin{equation}
 \begin{bmatrix}
  f_{\mathrm{aug}}(\theta_{\mathrm{aug}},\hat x_k,u_k)\\
  g_{\mathrm{aug}}(\theta_{\mathrm{aug}},\hat x_k,u_k)
 \end{bmatrix}
 =\phi_{\mathrm{aug}}(\theta_{\mathrm{aug}},\hat x_k,u_k),
 \label{eq:aug_mlp}
\end{equation}
where $f_{\mathrm{aug}}\in\R^6$, $g_{\mathrm{aug}}\in\R^{n_{x_a}}$, and
$\theta_{\mathrm{aug}}$ collects the weights and biases of
$\phi_{\mathrm{aug}}$.
$f_{\mathrm{aug}}$ corrects every physical row, the positions included, as
in the dynamic augmentation of~\cite[Sec.~2]{hoekstra2025lfr}, where the
learning component acts on the entire model state.
The correction is added to the result of the RK4 step, so
$f_{\mathrm{aug}}$ is a discrete-time correction, whereas $f_{\mathrm{base}}$
integrates the continuous-time physics.

\begin{figure}[t]
 \centering
 \resizebox{\columnwidth}{!}{\input{figures/tex/augmentation_structure}}
 \caption{Augmented model. The physics step $f_{\mathrm{base}}$ and the
 network $\phi_{\mathrm{aug}}$ act in parallel on $\hat x_k$.
 $\phi_{\mathrm{aug}}$ supplies $f_{\mathrm{aug}}$ and $g_{\mathrm{aug}}$.
 The output map is not learned ($h_{\mathrm{aug}}=0$). The encoder
 $\psi$ sets $\hat x$ at each window start. During training, $u_k$ is the
 closed-loop input $\hat u_k$.}
 \label{fig:aug_structure}
\end{figure}

The output map $h_{\mathrm{base}}$ is the kinematic map
$\begin{bmatrix}P^\top & 0\end{bmatrix}$ of \eqref{eq:Ptransform}, so the
measured positions are read from the physical state alone
(Fig.~\ref{fig:aug_structure}).
Within one step, no learned function feeds the input of another, so the
interconnection graph is acyclic.
Because $M(Y)$ is invertible for every admissible $Y$
(Section~\ref{sec:lfr}), $f_{\mathrm{base}}$ is twice continuously
differentiable.
With tanh activations, the model is therefore well
posed~\cite[Thm.~9]{hoekstra2026lfr}.

The additional states have no assigned physical meaning.
In particular, they are not identified with the states of the absorber that
the simulated system of Section~\ref{sec:setup} contains.
They affect the output only through $f_{\mathrm{aug}}$, one step later.

Training starts from the baseline.
The output layer of $\phi_{\mathrm{aug}}$ is initialised at zero, which gives
\begin{equation}
 \phi_{\mathrm{aug}}(\theta^0_{\mathrm{aug}},\hat x,u)=0
 \quad\text{for all }(\hat x,u),
 \label{eq:aug_zero_init}
\end{equation}
where $\theta^0_{\mathrm{aug}}$ denotes the initial network parameters.
The augmented model then reproduces the baseline from the same physical
initial state, as required by~\cite[Eq.~(29)]{hoekstra2026lfr}.

\subsection{Encoder and initial state}\label{sec:aug_encoder}
Each training window simulates \eqref{eq:aug_transition} from an initial
state, but only the positions are measured.
Following Beintema et al.~\cite{beintema2023subnet}, an encoder estimates
this state from the input and output history,
\begin{equation}
 \hat x_{\tau|\tau}
 =\psi_\eta(\xi_\tau),\qquad
 \xi_\tau=\col\bigl(y_{\tau-n:\tau},\,u_{\tau-n:\tau}\bigr),
 \label{eq:aug_encoder}
\end{equation}
where $\tau$ is the window start and $n$ the length of the history.
The history includes the current sample, as the reconstructability map
does~\cite[Eq.~(15)]{hoekstra2026encoder}.
The encoder estimates the full state, the measured positions included.

The encoder is a linear map with a nonlinear
correction~\cite[Eq.~(8)]{hoekstra2026encoder},
\begin{equation}
 \psi_\eta(\xi_\tau)
 =\begin{bmatrix}W^b\\ W^a\end{bmatrix}\xi_\tau+\tilde\psi(\xi_\tau),
 \label{eq:aug_enc_lin}
\end{equation}
where $W^b$ and $W^a$ give the physical and the additional states,
$\tilde\psi$ is an MLP, and $\eta$ collects $W^b$, $W^a$ and the parameters
of $\tilde\psi$.

The rows $W^b$ start from the reconstructability map of the linearised
baseline~\cite[Sec.~3.1]{hoekstra2026encoder}.
Linearising \eqref{eq:eom} about rest, with the nominal parameters of
Appendix~\ref{app:params}, gives
\begin{equation}
 A(Y)=\begin{bmatrix}0 & I\\ -M(Y)^{-1}K & -M(Y)^{-1}C\end{bmatrix},\qquad
 B(Y)=\begin{bmatrix}0\\ M(Y)^{-1}\end{bmatrix},
 \label{eq:aug_lin}
\end{equation}
with input $u$.
The map uses $A(0)$ and $B(0)$, discretised by zero-order hold over $T_s$ to
$(A_d,B_d)$, together with the matrix
$C_d=\begin{bmatrix}P^\top & 0\end{bmatrix}$ of $h_{\mathrm{base}}$.
For this model, $W^b$ is~\cite[Eqs.~(16), (17)]{hoekstra2026encoder}
\begin{equation}
 W^b=\begin{bmatrix}
  A_d^nO_n^{+} & \;r_n-A_d^nO_n^{+}T_n
 \end{bmatrix},
 \label{eq:aug_wb}
\end{equation}
where $O_n$ is the observability matrix of $(A_d,C_d)$ over the history $n$,
$T_n$ the Toeplitz matrix of input-to-output responses, $r_n$ the
input-to-state map over the same history, and the two blocks act on the
output and input parts of $\xi_\tau$.
All matrices are normalised as the signals.
Because all positions are measured, $(A_d,C_d)$ is observable, so $O_n$ has
full column rank and $O_n^{+}$ is a left inverse.
The linearisation sets only the starting value of $W^b$, because the encoder
is trained jointly with the model.

The rows $W^a$ have no baseline counterpart, so they are drawn randomly.
The correction $\tilde\psi$ starts at zero.
By \eqref{eq:aug_zero_init}, the encoded additional states do not affect the
output at initialisation.
They become active as the output layer of $\phi_{\mathrm{aug}}$ is trained.

\subsection{Closed-loop rollout and objective}\label{sec:aug_closed_loop}
The records are measured under feedback.
The model must also predict the machine under feedback, that is, its
response to a given reference, feedforward and known controller.
We therefore fit this closed-loop response, rather than the open-loop
simulation driven by the recorded input of~\cite[Eq.~(22)]{hoekstra2026lfr}.

The baseline gives a second reason.
$X$ and $Y$ have no stiffness in \eqref{eq:damping_stiffness_matrices}, so
$A(Y)$ in \eqref{eq:aug_lin} has two eigenvalues at the origin for every
$Y$.
In an open-loop rollout, a position error caused by a velocity error
therefore persists.
Inside the loop, the controller acts on such errors as it does on the
machine, provided the model in feedback with the controller is stable.

Kessels trains an augmented model with the known feedback controller inside
each truncated rollout and propagates gradients through
it~\cite[Eqs.~(5.12), (5.13c), (5.13d)]{kessels2025ai}.
We adopt this rollout.
As in~\cite[Sec.~5.3.2.3]{kessels2025ai}, the controller stays outside the
learned model, so that it can be replaced in evaluation.

The recorded loop is driven by a reference $r$ and an actuator feedforward
$u^{\mathrm{ff}}$ through a linear controller with realisation
$(A_{\mathrm c},B_{\mathrm c},C_{\mathrm c},D_{\mathrm c})$ and state $s$.
The model loop receives the same $r$, $u^{\mathrm{ff}}$ and controller.
It differs only in its output $\hat y$ (Fig.~\ref{fig:aug_closed_loop}):
\begin{equation}
\begin{aligned}
 s_{k+1}&=A_{\mathrm c} s_k+B_{\mathrm c}(r_k-y_k),\\
 u_k&=P\bigl(C_{\mathrm c} s_k+D_{\mathrm c}(r_k-y_k)+u^{\mathrm{ff}}_k\bigr),\\
 \hat s_{k+1}&=A_{\mathrm c} \hat s_k+B_{\mathrm c}(r_k-\hat y_k),\\
 \hat u_k&=P\bigl(C_{\mathrm c} \hat s_k+D_{\mathrm c}(r_k-\hat y_k)
   +u^{\mathrm{ff}}_k\bigr).
\end{aligned}
\label{eq:aug_direct_controller}
\end{equation}

Subtracting the recorded loop from the model loop removes $r$ and
$u^{\mathrm{ff}}$,
\begin{equation}
\begin{aligned}
 x^{\mathrm c}_{k+1}&=A_{\mathrm c} x^{\mathrm c}_k
   +B_{\mathrm c}(y_k-\hat y_k),\qquad x^{\mathrm c}_\tau=0,\\
 \hat u_k&=u_k+P\bigl(C_{\mathrm c} x^{\mathrm c}_k
   +D_{\mathrm c}(y_k-\hat y_k)\bigr),
\end{aligned}
\label{eq:aug_controller}
\end{equation}
where $x^{\mathrm c}=\hat s-s$ is the difference of the controller states.
The subtraction is exact when both loops apply the same linear controller,
not scheduled on the model state, with the same feedforward and without
saturation.
The model loop runs the controller at the model sample period $T_s$, which
differs from the recording rate.
Section~\ref{sec:setup} reports the effect on the loop.

\begin{figure}[t]
 \centering
 \includegraphics[width=\columnwidth]{closed_loop_same_reference.pdf}
 \caption{Recorded plant loop (top) and augmented-model loop (bottom),
 driven by one shared reference $r_k$ and feedforward $u_k^{\mathrm{ff}}$,
 with the same feedback controller. The model output is evaluated
 before the input advances its state to the next sample.}
 \label{fig:aug_closed_loop}
\end{figure}

The condition $x^{\mathrm c}_\tau=0$ gives the model's controller the state
of the recorded controller, $\hat s_\tau=s_\tau$.
Model error before $\tau$ is therefore not carried into the window.
Kessels reconstructs this state from the measured output, the reference and
the controller~\cite[Remark~5.4]{kessels2025ai}.
In \eqref{eq:aug_controller}, the recorded input already carries it, so no
reconstruction is needed.

Because $h_{\mathrm{base}}$ has no feedthrough, each sample is evaluated in a
fixed order: $\hat y_k$ from the state, then $\hat u_k$ from
\eqref{eq:aug_controller}, then both states advance, with $\hat u_k$ held
over the RK4 step.

The controller suppresses model error inside its bandwidth.
A good closed-loop fit therefore does not by itself show that the open-loop
plant is recovered.

Over window starts $\mathcal T$ and window length $n_f$, the parameters
$\vartheta=(\theta_{\mathrm{base}},\theta_{\mathrm{aug}},\eta)$ minimise the
normalised output error
\begin{equation}
 V(\vartheta)=\frac{1}{|\mathcal T|\,n_f}
 \sum_{\tau\in\mathcal T}\sum_{\ell=0}^{n_f-1}
 \norm{y_{\tau+\ell}-\hat y_{\tau+\ell|\tau}}_2^2
 \label{eq:aug_loss}
\end{equation}
subject to \eqref{eq:aug_transition} with $u_k=\hat u_k$,
\eqref{eq:aug_encoder} and \eqref{eq:aug_controller}.
This is the truncated objective of~\cite[Eq.~(22)]{hoekstra2026lfr}, with
each window simulated in closed loop as in~\cite[Eq.~(5.12)]{kessels2025ai}.
As in both, every sample of the window is scored.

The window length $n_f$ is set by the memory of the closed loop: the window
must span the time over which a force error of the model remains visible in
the output.
Validation uses a longer horizon.
Checkpoints are selected on a closed-loop simulation over complete
validation records, started once from the encoder, as
in~\cite[Eq.~(5.14)]{kessels2025ai}.
Section~\ref{sec:setup} gives $n_{x_a}$, the network size, the encoder
history $n$, the window length $n_f$ and the validation records.
